\section{Every-power inertia deficits in an alternating partner layer}
\label{sec:projected-inertia-power-deficit}

\paragraph{Status and shared provenance.}
D-TD47, reserved question Q2fc, is a separate derivation by
\texttt{/root/imprimitive}. Root and the author independently derived
the composite-order deficit estimate; root supplied its stronger
prime-order version and the infinite prime-triangle application.
The local normalizer interface was reconstructed jointly with
\texttt{/root/sixfolds}, whose D-TD46 treats the sharper
\((2,3,7)\) case. The argument below keeps the actual original
source and every surviving inertia power. The complete proof, source interfaces, checker and first successful
output have passed the author's full reading. Independent and root whole-proof
reviews also pass after the recorded comparison-scope repair; no separate origin is claimed for the shared steps.

\subsection*{An actual permutation quotient controls every inertia power}

Let \(Z\longrightarrow\mathbf P^1\) be a connected finite
\(Q\)-Galois cover with genuine generating product-one inertia
\(q_1,\ldots,q_s\) of exact orders \(d_i\ge2\). Put
\[
 \Delta=-2+\sum_{i=1}^s(1-1/d_i)>0,\qquad
 \kappa(d)=
 \begin{cases}
 d/(d-1),&d\text{ is prime},\\
 d,&d\text{ is composite},
 \end{cases}
 \qquad \kappa=\max_i\kappa(d_i).
 \tag{IPD.1}
\]
Suppose \(Q\) has an actual transitive permutation action on
\(\Omega\), of degree \(D\), and write \(y_i\) for the image
of \(q_i\). The action may have a kernel. Define
\[
 \varepsilon_i=(1-1/d_i)D-\operatorname{ind}_\Omega(y_i),
 \qquad
 \theta=1-\kappa\Delta-2\kappa/D.
 \tag{IPD.2}
\]

\begin{lemma}[D-TD47, every-power deficit]
One has \(\varepsilon_i\ge0\) and
\[
 \sum_i\varepsilon_i\le\Delta D+2.
 \tag{IPD.3}
\]
For every \(i\) and every integer \(k\) for which
\(q_i^k\ne1\),
\[
 \operatorname{supp}_\Omega(y_i^k)
 \ge D-\kappa(d_i)\varepsilon_i
 \ge(1-\kappa(d_i)\Delta)D-2\kappa(d_i)
 \ge\theta D.
 \tag{IPD.4}
\]
In particular a positive \(\theta\) rules out an identity
permutation image of any such nonidentity inertia power.
\end{lemma}

\begin{proof}
Every cycle length \(l\) of \(y_i\) divides \(d_i\), even
when the order of \(y_i\) is a proper divisor of \(d_i\).
Writing the index as the sum of \(l-1\) over the cycles gives
\[
 \varepsilon_i=\sum_{\text{cycles of }y_i}(1-l/d_i)
 =\sum_{\omega\in\Omega}(1/l_\omega-1/d_i)\ge0.
 \tag{IPD.5}
\]
The connected quotient cover of degree \(D\) has some genus
\(g_\Omega\ge0\). Its actual Riemann--Hurwitz equation is
\(\sum_i\operatorname{ind}(y_i)=2D-2+2g_\Omega\).
Subtracting this from \((2+\Delta)D\) proves
\(\sum_i\varepsilon_i=\Delta D+2-2g_\Omega\), hence (IPD.3).

Replace \(k\) by its residue modulo \(d_i\), between one
and \(d_i-1\). A letter fixed by \(y_i^k\) lies on a cycle
whose length \(l\) divides both \(d_i\) and \(k\), so
\(l\) is a proper divisor of \(d_i\). If \(d_i\) is
composite, \(l\le d_i/2\) and
\(1/l-1/d_i\ge1/d_i=1/\kappa(d_i)\).
If \(d_i\) is prime, necessarily \(l=1\), giving instead
\(1/l-1/d_i=(d_i-1)/d_i=1/\kappa(d_i)\).
Thus each fixed letter contributes at least \(1/\kappa(d_i)\)
to (IPD.5). All other contributions are nonnegative. Consequently
the number of fixed letters is at most
\(\kappa(d_i)\varepsilon_i\). Use (IPD.3), then
\(\kappa(d_i)\le\kappa\), to obtain every inequality in
(IPD.4). This proof never substitutes the projected order for
the exact order in the actual cover.
\end{proof}

\subsection*{The complete alternating-layer normalizer interface}

Suppose that a faithful finite complex representation \(U\) of
\(Q\), of degree \(b\), has a stable whole tensor layer
\[
 U=(W_1\otimes\cdots\otimes W_t)\otimes R,
 \qquad m\ge8,\qquad D=mt.
 \tag{IPD.6}
\]
All factors and the entire complement \(R\) are retained.
The actual projective image on the displayed layer contains the
full normal core \(A_m^t\), represented through its perfect
central lifts on the irreducible spaces \(W_j\), and permutes
the component algebras transitively. This is the actual coarse
nonabelian extraction layer of D-TD12; ordinary extendibility,
standard labels and trivial central action are not assumed.
Conjugation gives the actual natural action on \(mt\) letters
through \(\operatorname{Aut}(A_m^t)=S_m\wr S_t\).
It is transitive, because the full inner core is transitive
inside each block and the component action is transitive.

\begin{lemma}[D-TD47, restricted whole-partner floor]
If \(\theta>0\) in (IPD.2) for this actual natural action,
then every nonidentity inertia power in the stated tuple satisfies
\[
 \frac{\nu_U(q_i^k)}{b}\ge\frac{\theta}{24},\qquad
 \nu_U(T)=b-\max_z\dim\ker(U(T)-zI).
 \tag{IPD.7}
\]
This is a floor on these powers, not a claim about every
nonidentity element of \(Q\).
\end{lemma}

\begin{proof}
First justify the support-sensitive estimate on a fixed component.
Its represented perfect central lift \(E\) of \(A_m\) acts
irreducibly and nontrivially on \(W_j\). A compatible local
operator \(T\) normalizes \(E\) projectively. Perfectness
promotes this to honest normalization: taking derived groups in
\(TET^{-1}\subset\mathbf C^*E\) gives containment in \(E\),
and applying the same argument to \(T^{-1}\) gives equality.
The projective order of \(T\) is finite. Rescale it so that
\(T\) has finite order; \(G=\langle E,T\rangle\) is then
finite, since \(E\) is finite and normalized by \(T\).

Conjugation on \(E/Z(E)\) gives a map from \(G\) into
\(\operatorname{Aut}(A_m)=S_m\), whose image contains \(A_m\).
An element of its kernel conjugates each \(e\in E\) to
\(c(e)e\). The scalar function \(c\) is a character of the
perfect group \(E\), so is trivial. Schur's lemma makes the
kernel scalar, and hence central in \(G\). If the natural
image \(y\) of \(T\) is even, choose an element of \(E\)
with the same image; their ratio is scalar. The even case of
D-TD31 applies. If \(y\) is odd, the map onto \(S_m\) and
this actual representation supply every hypothesis of the
precise odd-extension variant of D-TD31. Therefore in either
case, for \(y\ne1\),
\[
 \frac{\nu_{W_j}(T)}{\dim W_j}
       \ge\frac{\operatorname{supp}_m(y)}{24m}.
 \tag{IPD.8}
\]
No actual lift of \(E\) to the original coefficient group or
to \(Q\) is asserted; the finite group \(G\) was constructed
in this actual local operator space. Scalar normalization does
not change projective spectral rank.

Now consider \(q_i^k\ne1\). If it moves a component algebra,
one moved cycle of length at least two has the accepted quarter
support bound of HD-FT18: in an eigenbasis for its finite-order
cycle product, all nonconstant words have rotation orbits of
length at least two. Each weighted orbit contributes at most
one vector to a prescribed eigenspace. Its dimension is at
least two in each factor, giving normalized support at least
\(1/4\). Spectral convolution with the rest of the layer and
\(R\) preserves this bound.

Otherwise every component is fixed. By (IPD.4), the sum of
the natural supports on its \(t\) blocks is at least
\(\theta mt\), so at least one component moves at least
\(\theta m\) letters. Apply (IPD.8) and again convolve with
every other factor, including \(R\). This gives (IPD.7).
Finally \(0<\theta<1\) because \(\Delta>0\), so the
quarter bound in the moved case is stronger than \(\theta/24\).
Neither case singles out a stable component quotient from
a permuted layer.
\end{proof}

\subsection*{A bound on the original standard degree}

Use the actual standard-source setting of HD-SS01--02:
\[
 \begin{gathered}
 C\longrightarrow\mathbf P^1\text{ a connected finite }J\text{-cover},
 \quad V=(\operatorname{Std}_n\circ x)\otimes U,
 \quad x(J)\in\{A_n,S_n\},\quad n\ge5,\\
 b=\dim U\ge2,\quad \mu_V=6,\quad
 K_0=\ker U\text{ transitive on the natural }n\text{ letters},
 \quad Q=U(J),\quad \mu=\mu_U>0.
 \end{gathered}
 \tag{IPD.9}
\]
Here \(V\) is faithful and absolutely irreducible, and both
factors are actual finite linear representations. Give the
actual \(Q\)-cover \(C_U=C/K_0\to\mathbf P^1\) its genuine
orders \(d_i\). Suppose the same entire \(U\) contains a layer
(IPD.6), with \(\theta>0\). Put
\[
 \omega=\max_i(1-1/d_i),\qquad c=24\omega/\theta.
 \tag{IPD.10}
\]

\begin{theorem}[D-TD47, mixed-order source bound]
In this setting,
\[
 n\Delta\le c\left(2+\frac{6+\mu}{b}\right)-2
 \le c\left(4+\frac6b+
                  \frac{4+12/b}{n-2}\right)-2.
 \tag{IPD.11}
\]
In particular, for \(n\ge12\),
\[
                         n\Delta\le8c-2.
 \tag{IPD.12}
\]
The inequalities allow every genus of the natural quotient
\(X=C/x^{-1}((S_n)_1)\). No source budget is imposed on \(U\)
as a replacement for the original representation \(V\).
\end{theorem}

\begin{proof}
The irreducible \(U\) has degree at least two. Quotient pullback
identifies its multiplicity in \(H^1(C)\) with its multiplicity
in \(H^1(C_U)\), since \(K_0\) acts trivially on \(U\).
The latter is positive by hypothesis, so \(g(C_U)\ge2\): genus zero has no cohomology,
and every irreducible constituent of the cohomology of a
genus-one curve has dimension one. Thus the \(\Delta\) of
(IPD.1) is indeed positive.

Let \(J_1=x^{-1}((S_n)_1)\), \(N=J_1\cap K_0\),
\(X=C/J_1\), \(Y=C/N\), and \(h=g(X)\).
Kernel transitivity gives \(J=K_0J_1\),
\(J_1/N\simeq Q\), \([K_0:N]=n\), and
\(U(J_1)=U(J)\). The cover \(Y\to X\) is the connected
normalization of the actual base change of \(C_U\) along \(X\).
The map \(Y\to C_U\) has degree \(n\).

By the complete local calculation in HD-SS11, its inertia
operators are conjugates of powers \(q_i^l\), where \(l\)
is a natural cycle length at the original branch point. Their
orders are \(d_i/\gcd(d_i,l)\); trivial powers are omitted
as genuine inertia. No original lift order is replaced by
the projected order in making this assertion. If a descended
order is \(o>1\), its weight is at most \(\omega\), while
(IPD.7) bounds its actual moving rank below by \(b\theta/24\).
Writing \(D'\) for the complete descended moving-rank sum gives
\[
 \sum_{\text{descended genuine marks}}(1-1/o)
 \le\frac cb D',\qquad
                  D'=6+\mu-2b(h-1).
 \tag{IPD.13}
\]
The second identity is exactly HD-SS02, with its extra
\(U\)-cohomology term retained.

Riemann--Hurwitz for the two actual maps from \(Y\) now gives
\[
 n\Delta\le\frac{2g(Y)-2}{|Q|}
 \le2h-2+\frac cb\{6+\mu-2b(h-1)\}
 =\frac{c(6+\mu)}b+(2-2c)(h-1).
 \tag{IPD.14}
\]
Since \(\omega\ge1/2\) and \(0<\theta<1\), one has
\(c>12>1\). The right side is therefore largest at \(h=0\),
which proves the first inequality in (IPD.11).
The accepted all-inertia bound (SS.17),
\(\mu\le2b+(4b+12)/(n-2)\), proves the second.
For \(b\ge2\) and \(n\ge12\), its parenthesis is at most
\(7+10/(n-2)\le8\), giving (IPD.12).
Every step uses the original source, the exact descended
cohomology identity and actual quotient covers.
\end{proof}

\subsection*{All prime signatures \((2,3,p)\), not only \((2,3,7)\)}

\begin{corollary}[D-TD47, infinite prime-signature exclusion]
Suppose (IPD.9) holds, the actual \(Q\)-cover has signature
\((2,3,p)\) for a prime \(p\ge7\), and the entire partner
contains a layer (IPD.6) with \(D=mt\ge24\). Then
\[
 b\ge3,\qquad
 n\le1332+\frac{6648}{p-6}.
 \tag{IPD.15}
\]
In particular, for every prime \(p\ge11\), one has \(n\le2661\).
The same calculation gives \(n\le7980\) for \(p=7\);
the separate D-TD46 bound \(4596\) is stronger when its additional
exact imaginary-quadratic \(K\)-model and \(D\ge27\) hypotheses hold.
The new infinite family here is the range \(p\ge11\).
\end{corollary}

\begin{proof}
Here
\[
 \Delta=\frac16-\frac1p,\qquad \kappa=2,\qquad
 \theta=1-2\Delta-4/D>\frac12,
 \qquad \omega=1-1/p.
 \tag{IPD.16}
\]
The inequality follows from \(D\ge24\) and
\(\Delta<1/6\). Thus every genuine descended inertia has
normalized projective support at least \(1/48\). Its actual
order is one of \(2,3,p\), by the normalized base-change
calculation, and its orbifold weight is at most \(1-1/p\).

The \((2,3,p)\) triangle group has trivial abelianization.
Indeed in an abelian quotient its three generators satisfy
\(2u=3v=pw=0\) and \(u+v+w=0\); multiplication by six,
then coprimality of \(p\) and six, forces \(w=0\), and
the coprime orders two and three force \(u=v=0\).
Its actual quotient \(Q\) is therefore perfect, so
\(\det U=1\). The order-two inertia cannot act as a scalar
under faithful \(U\). Otherwise it is central in \(Q\),
and product one expresses the third generator through the
first two, making \(Q\) abelian. This contradicts its
perfectness and its genuine inertia orders.

Consequently that involution has both eigenvalues \(1,-1\).
It fixes at least one vector. Euler on the actual three-mark
\(Q\)-sphere cover gives
\[
 \mu=-2b+\sum_{i=1}^3\operatorname{rank}(1-U(q_i))
                      \le b-1.
 \tag{IPD.17}
\]
If \(b=2\), a determinant-one nonidentity involution would
instead be \(-I\), which is scalar. Hence \(b\ge3\).

Repeat (IPD.14) with the valid coarser coefficient
\(48(1-1/p)\) in place of \(c\). Using \(\mu\le b-1\)
and \(b\ge3\), the maximum at \(h=0\) gives
\[
 n\left(\frac16-\frac1p\right)
 \le48(1-1/p)\left(3+\frac5b\right)-2
 \le222-\frac{224}{p}.
 \tag{IPD.18}
\]
Solving this exact rational inequality proves (IPD.15).
For \(p\ge11\) the right side in (IPD.15) is at most
\(1332+6648/5=13308/5\), so the integer \(n\) is at most
2661. For \(p=7\) it equals 7980. These are necessary source
bounds, not constructions or reverse implications.
\end{proof}

\paragraph{Application and remaining boundary.}
The local and source theorems impose no standard, ordinary,
extendible or singleton condition on the selected alternating
layer of \(U\). Its genuine central lifts and component
permutations are all retained. They introduce no new external
classification theorem: the precise local input is accepted
D-TD31, and all cover and cohomology inputs are the accepted
HD-SS01--02/06/11 interfaces.

In a minimal exact-\(K\), rank-nine, naturally polarized
maximal-standard source, the separately proved kernel and
post-vanishing descent conclusions supply transitivity and
\(\mu>0\). Once the complete-complement bound D-TD44 supplies
an alternating extraction layer of degree at least
\(\lceil(n+120)/108\rceil\), a sufficiently large such source
has \(m\ge8\) and \(mt\ge24\). Thus its actual
\(Q\)-signature cannot be \((2,3,p)\) for any prime
\(p\ge11\); D-TD46 separately disposes of \(p=7\).
This last application is conditional on those separately stated
inputs at their recorded review status. The standalone theorem
above does not depend on D-TD44 or D-TD46.

When \(\theta\le0\), (IPD.4) remains a valid inequality but
does not supply the positive whole-partner floor used here.
No global bound on the remaining original degree, partner
degree or inducing image follows, and no sphere budget has
been iterated on \(U\).
